\section{Related Work and Taxonomy}
\label{sec:related}

Table~\ref{tab:taxonomy} summarizes twenty peer-reviewed systems and performance-modeling papers from ACM, USENIX, and other peer-reviewed venues. The selection is a targeted search and not a PRISMA-2020 systematic review, so it does not claim exhaustive coverage. The table records each work's scope, method, and the limitation that matters for multi-tier thread-pool analysis. This section groups the same works into four threads and states what each one leaves open.

\subsection{Tail-latency amplification in microservice call graphs}

Dean and Barroso~\cite{dean2013tailatscale} showed that parallel fan-out amplifies tail latency through order statistics, because a parent request completes only when its slowest branch returns. Their argument is qualitative and does not model how a branch's queueing delay propagates through the holding time of its parent. Gan et al.~\cite{gan2019deathstarbench} introduced DeathStarBench, an end-to-end benchmark of microservice applications, and reported that RPC paths in these applications have substantially higher tail-to-median ratios than monolithic workloads. The benchmark is the source of the two sub-graphs used in our Phase~1 and Phase~3 campaign. Sriraman and Wenisch~\cite{sriraman2018musuite} provide a complementary microservice benchmark suite that reports tail latency for its services. Redundant requests address the same tail from the other side: cloning reduces stragglers at the cost of extra work~\cite{ananthanarayanan2013clones}, and the tail-at-scale argument favors hedged reads for idempotent calls~\cite{dean2013tailatscale}. Tail-latency guarantees for shared resources are studied by PriorityMeister~\cite{zhu2014prioritymeister} and SNC-Meister~\cite{zhu2016snc}, which configure priorities and admission against tail SLOs in shared storage and networks rather than in microservice graphs.

Luo et al.~\cite{luo2021characterizing} characterized microservice dependencies and performance in production traces from Alibaba clusters. Their analysis finds that a small number of heavily shared downstream services dominate the dependency graphs, which is the structure behind the diamond in-cast graphs we study. Zhou et al.~\cite{zhou2021fault} studied fault analysis and debugging of microservice systems through an industrial survey, a benchmark system, and an empirical study. Their focus is fault diagnosis rather than the latency that builds up before a fault is declared.

The learning-based line of work uses historical traces to act on microservice performance. Seer~\cite{gan2019seer}, Sage~\cite{gan2021sage}, Sinan~\cite{zhang2021sinan}, and FIRM~\cite{qiu2020firm} diagnose QoS violations or adjust resource allocations in a closed loop. These systems treat each tier as a black-box compute node. None of them isolates how synchronous thread holding couples the saturation of an upstream pool to queueing at a downstream tier that is still below its CPU limit. That coupling is the subject of the present model. Cluster schedulers that account for interference and QoS act on the same problem at the placement level, as in Paragon~\cite{delimitrou2013paragon} and Quasar~\cite{delimitrou2014quasar}, but they do not model the holding of threads along a call path.

\subsection{Concurrency runtimes, thread-pool saturation, and metastability}

When an upstream handler blocks a worker thread while it awaits downstream responses, queueing delay at a leaf tier directly inflates the upstream holding time. Bronson et al.~\cite{bronson2021metastable} and Huang et al.~\cite{huang2022metastable} described metastable failures, in which overload feedback keeps a system degraded after the trigger has passed. Both works describe how retries and load feedback can hold a system in a degraded state, but they do not measure how that feedback depends on the depth and width of a call graph.

Staged event-driven designs made the thread pool itself a control point for overload~\cite{welsh2001seda}. Recent microsecond-scale schedulers preempt or steer requests to contain tails~\cite{kaffes2019shinjuku,ousterhout2019shenango}, and they confirm that scheduling rather than raw capacity often sets the tail. The measurement methodology of this paper rests on a result about load generation. Schroeder et al.~\cite{schroeder2006open} showed that closed-loop generators understate the saturation knee of a server, because they throttle arrivals whenever the server slows down. Open-loop schedules with coordinated-omission correction avoid this bias, and our generator measures each request's latency from its intended send time. The statistical tools used to locate the saturation horizon and to fit tails are classical. Whitt~\cite{whitt1993ggm} gives approximations for $GI/G/m$ stations that underlie the waiting-time term of our model. Kneedle~\cite{satopaa2011kneedle} locates the knee of a curve from its curvature, bootstrap resampling~\cite{efron1979bootstrap} gives interval estimates without distributional assumptions, and peaks-over-threshold extreme value theory~\cite{pickands1975extreme,davison1990exceedances} gives the generalized Pareto model for excesses above a high threshold. To our knowledge these tools have not been combined to test a multi-tier holding-time recursion across synchronous, asynchronous, and hybrid runtimes.

\subsection{Service mesh sidecars, serialization, and observability overheads}

Service meshes interpose user-space proxies and telemetry hooks on every RPC edge. Karandikar et al.~\cite{karandikar2021protoacc} designed a hardware accelerator for Protocol Buffers after finding that serialization is a large share of cost in datacenter workloads. Verma et al.~\cite{verma2015borg} documented resource interference and scheduling behavior in Google's Borg cluster manager, which is the environment in which such sidecars are deployed. Zhu et al.~\cite{zhu2023sidecar} dissected the overheads of service mesh sidecars, and that study is the closest published work to the sidecar component of our Phase~3 experiment. Distributed tracing systems address causal context propagation. Pivot Tracing~\cite{mace2015pivot} propagates baggage along request paths, and Canopy~\cite{kaldor2017canopy} describes tracing at production scale.

Controlled experiments that separate per-hop serialization from sidecar socket forwarding across call-graph depths are scarce. The same holds for experiments that bound how span sampling inflates the tail near thread-pool saturation. Phase~3 and Phase~4 of this study address those two gaps directly, with an emulated sidecar and an instrumented chain, and they should be read as mechanism-level measurements rather than as production figures.

\subsection{Overload control, circuit breaking, and deadline-aware shedding}

Admission control and circuit breaking are the standard defenses against overload collapse. DAGOR~\cite{zhou2018dagor} sheds requests by priority, keyed on queueing delay at each service, and it was deployed for WeChat's microservices. Breakwater~\cite{cho2020breakwater} uses credit-based server-side admission for microsecond-scale RPCs. Both controllers protect a single tier. A local controller that ignores the remaining end-to-end budget can admit a request at an upstream hop that will expire before its critical path completes. Static circuit breakers and naive retries can also amplify a capacity fault, because retries add load to a tier that is already saturated. DAC-Shed combines hop-by-hop deadline-budget shedding with a latency-gradient concurrency limit adapted from TCP Vegas~\cite{brakmo1994vegas}. The gradient rule compares each 50\,ms window with the smallest of the last twenty windows, which is the delay-based signal that Vegas uses for congestion. Scheduling that favors short work under overload has a long history~\cite{schroeder2006web}, and the controller here applies a delay signal to admission rather than to request order.

\subsection{Positioning}

The four threads share a gap. Trace studies characterize the graphs, learning systems act on them, and admission controllers protect single tiers. None of the verified works, within the search scope above, measures coupled thread holding across a bounded call graph under controlled open-loop load with coordinated-omission correction, and none evaluates a deadline-budget controller against a static breaker and naive retries under the same injected faults. The contributions in Section~\ref{sec:intro} are aimed at that gap.

\begin{table*}[t]
\centering
\footnotesize
\setlength{\tabcolsep}{3pt}
\caption{Curated taxonomy of 20 peer-reviewed systems papers (ACM, USENIX, and other peer-reviewed venues). Every record was checked against Crossref or the USENIX presentation page. The selection is targeted, not a PRISMA-2020 review, and the limitations describe scope rather than quality.}
\label{tab:taxonomy}
\resizebox{\linewidth}{!}{%
\begin{tabular}{@{}p{2.6cm}lp{3.2cm}p{3.3cm}p{5.2cm}@{}}
\toprule
\textbf{Work} & \textbf{Venue} & \textbf{Scope} & \textbf{Method} & \textbf{Limitation for multi-tier thread-pool analysis} \\
\midrule
Dean and Barroso \cite{dean2013tailatscale} & CACM 2013 & Fan-out request paths & Tail-tolerant techniques & Tail amplification argued qualitatively; no coupled multi-tier queueing model \\
Gan et al.\ DeathStarBench \cite{gan2019deathstarbench} & ASPLOS 2019 & End-to-end microservice applications & Hardware and software characterization & Characterizes applications; does not isolate the concurrency runtime or DAG depth \\
Gan et al.\ Seer \cite{gan2019seer} & ASPLOS 2019 & Cloud microservice traces & Learned QoS-violation prediction & Prediction-oriented; no analysis of thread holding times \\
Gan et al.\ Sage \cite{gan2021sage} & ASPLOS 2021 & Microservice dependency graphs & ML culprit localization & Diagnosis after the fact; no admission control \\
Zhang et al.\ Sinan \cite{zhang2021sinan} & ASPLOS 2021 & Microservice resource allocation & ML-predicted end-to-end tail latency & Allocation-level control; thread-pool saturation below CPU capacity not modeled \\
Qiu et al.\ FIRM \cite{qiu2020firm} & OSDI 2020 & Multi-tier microservices & Learned SLO-driven resource reallocation & Container-level reallocation; no deadline-aware shedding \\
Luo et al.\ \cite{luo2021characterizing} & SoCC 2021 & Production microservice traces & Dependency and runtime characterization & Observational; no controlled concurrency-model experiments \\
Zhou et al.\ \cite{zhou2021fault} & IEEE TSE 2021 & Microservice fault analysis & Industrial survey and benchmark system & Focus on fault debugging rather than latency saturation \\
Zhou et al.\ DAGOR \cite{zhou2018dagor} & SoCC 2018 & WeChat microservices & Priority-based overload control & Evaluated on one production system; no DAG-coupled analysis \\
Huang et al.\ \cite{huang2022metastable} & OSDI 2022 & Incident reports across organizations & Metastable failure taxonomy & Qualitative incident study; no controlled latency measurements \\
Bronson et al.\ \cite{bronson2021metastable} & HotOS 2021 & Distributed systems & Metastable failure model & Conceptual; not evaluated on DAG topologies \\
Cho et al.\ Breakwater \cite{cho2020breakwater} & OSDI 2020 & Microsecond-scale RPC servers & Credit-based server-side overload control & Server-side admission; multi-tier coupling is not the focus \\
Brakmo et al.\ TCP Vegas \cite{brakmo1994vegas} & SIGCOMM 1994 & TCP congestion control & RTT-based delay-gradient control & Transport-layer method; the gradient controller here is adapted from it \\
Satop\"a\"a et al.\ Kneedle \cite{satopaa2011kneedle} & ICDCS-W 2011 & System behavior curves & Knee-point detection & General detector; no bootstrap intervals \\
Schroeder et al.\ \cite{schroeder2006open} & NSDI 2006 & Workload generation & Open versus closed system models & Establishes the open-loop requirement; not microservice-specific \\
Karandikar et al.\ ProtoAcc \cite{karandikar2021protoacc} & MICRO 2021 & Protocol Buffers (de)serialization & Hardware accelerator & Hardware offload; general-purpose CPU serialization tax not characterized here \\
Verma et al.\ Borg \cite{verma2015borg} & EuroSys 2015 & Cluster management & Large-scale scheduling & Scheduling, not request-path latency \\
Mace et al.\ Pivot Tracing \cite{mace2015pivot} & SOSP 2015 & Distributed monitoring & Dynamic causal monitoring & Monitoring design; tail-percentile observer effect not measured \\
Kaldor et al.\ Canopy \cite{kaldor2017canopy} & SOSP 2017 & Production tracing & End-to-end performance tracing & Tracing infrastructure; not a controlled observer-effect study \\
Zhu et al.\ \cite{zhu2023sidecar} & SoCC 2023 & Service mesh sidecars & Overhead decomposition & Sidecar overhead in production-style proxies; not DAG-coupled queueing \\
\bottomrule
\end{tabular}%
}
\end{table*}

